\begin{frame}{Method 2: General Access Structures \cite{beimel2011secret}}
	\textbf{\textcolor{blue}{Why general access structures?}}
	\begin{itemize}
		\item Real policies are rarely "any \(t\) out of \(n\)".
		\item Example: \((\text{CEO} \wedge \text{CFO}) \vee (\text{3 of 5 board members})\).
		\item Secret sharing can realize any monotone \(\mathcal{A}\), but \textbf{efficiency} is the main challenge.
	\end{itemize}

	\textbf{\textcolor{blue}{Key idea}}
	\begin{itemize}
		\item Describe \(\mathcal{A}\) via a monotone formula / span program.
		\item Implement AND/OR composition by combining smaller schemes along the formula structure.
	\end{itemize}
\end{frame}

\begin{frame}{Current Methods 3: LSSS and CRT-based Schemes}
	\begin{itemize}
		\item Many efficient constructions are \textbf{linear}: shares are linear combinations of secret + randomness over a field.
		\item A common abstraction uses:
			\begin{itemize}
				\item a matrix \(M\) over \(\mathbb{F}_q\),
				\item a labeling \(\rho\) mapping each row to a participant.
			\end{itemize}
		\item Authorized sets can reconstruct because their rows "span" the secret component; unauthorized sets cannot (perfect privacy) \cite{beimel2011secret}.
		\item Another family: CRT-based secret sharing (shares are residues modulo carefully chosen integers) \cite{iftene2007general}.
	\end{itemize}
\end{frame}
